% Einheit 2 von 5 – Lineare Funktion f(x) = m·x + n (bank/lineare-funktionen/e2.jsonl, zusammenbau v0.1)
%% TODO Zweigzeile Teil 1: was hier gelernt wird (Satz in Schülersprache aus dem Einheitstitel) – Einheit 2
\einheitenkopf[e2]{Einheit 2 von 5 · Lineare Funktion f(x) = m·x + n}
\zweigzeile{ab Kl. 8 · P10 oft}
\setcounter{aufgabe}{11}

%% TODO Titel: Ich-kann-Satz fehlt in der Bank (Platzhalter: Kettenname) – Nr. 12
%% TODO Anweisung: ein Satz über der ersten Teilaufgabe fehlt in der Bank – Nr. 12
\begin{aufgabe}{m und n in der Gleichung markieren}
\begin{teile}
\teil m und n? Unterstreiche m, kreise n ein: $f(x) = 3x - 5$ m = \leerfeld , n = \leerfeld
\end{teile}
\end{aufgabe}

%% TODO Titel: Ich-kann-Satz fehlt in der Bank (Platzhalter: Kettenname) – Nr. 13
%% TODO Anweisung: ein Satz über der ersten Teilaufgabe fehlt in der Bank – Nr. 13
\begin{aufgabe}{Steigungsrichtung ohne Zeichnen}
\begin{teile}
\teil Steigt oder fällt die Gerade $f(x) = 3x - 4$? Kreuze an. \\ \kreuz{steigt} \\ \kreuz{fällt}
\end{teile}
\end{aufgabe}

%% TODO Titel: Ich-kann-Satz fehlt in der Bank (Platzhalter: Kettenname) – Nr. 14
%% TODO Anweisung: ein Satz über der ersten Teilaufgabe fehlt in der Bank – Nr. 14
\begin{aufgabe}{Steigungsdreieck lesen}
\begin{teile}
\teil Einen Schritt nach rechts – wie viel nach oben oder unten? Lies am Steigungsdreieck ab.

\begin{ksys}[xmin=-4,xmax=4,ymin=-5,ymax=5,ablesen] \gerade{3}{-1}{g} \steigungsdreieck{0}{-1}{3} \end{ksys}
\leerfeld nach \leerfeld \leerfeld
\end{teile}
\end{aufgabe}

%% TODO \verfahren{…}: Verfahrensüberschrift in Sachsprache fehlt in der Bank (2.3 b)
%% TODO Titel: Ich-kann-Satz fehlt in der Bank (Platzhalter: Kettenname) – Nr. 15
%% TODO Anweisung: ein Satz über der ersten Teilaufgabe fehlt in der Bank – Nr. 15
\begin{aufgabe}{Graph zeichnen}
%% TODO \rechenplatz{n}: Schrittzahl des Grundfalls fehlt in der Bank (nur bei fester Schreibform, 2.3 b)
\begin{teile}
\teil Wertetabelle: Fülle die Tabelle zu $f(x) = 3x - 2$ aus.

\wertetabelle{x}{f(x)}{-1,0,1,2}
\teil Punkte eintragen: Trage die Wertepaare $(-2|-2)$, $(-1|1)$, $(0|4)$, $(1|7)$ der Funktion $f(x) = 3x + 4$ ins Koordinatensystem ein.

\begin{ksys}[xmin=-4,xmax=4,ymin=-4,ymax=8] \end{ksys}
\teil Gerade durch Punkte: Berechne zwei Punkte von $f(x) = 4x + 2$, trage sie ein und zeichne die Gerade durch sie.

\begin{ksys}[xmin=-4,xmax=4,ymin=-4,ymax=7] \end{ksys}
\teil Mit n und Steigungsdreieck: Zeichne die Gerade zu $f(x) = 3x + 5$: Markiere n auf der y-Achse, gehe 1 nach rechts und m nach oben.

\begin{ksys}[xmin=-4,xmax=4,ymin=-4,ymax=9] \end{ksys}
\teil Mit n und Steigungsdreieck: Zeichne die Gerade zu $f(x) = -3x + 4$.

\begin{ksys}[xmin=-4,xmax=4,ymin=-4,ymax=5] \end{ksys}
\teil Mit n und Steigungsdreieck: Zeichne die Gerade zu $f(x) = \frac{1}{2}x + 3$.

\begin{ksys}[xmin=-4,xmax=4,ymin=-4,ymax=5] \end{ksys}
\teil Mit n und Steigungsdreieck: Zeichne die Gerade zu $f(x) = 2x - 4$.

\begin{ksys}[xmin=-4,xmax=4,ymin=-5,ymax=4] \end{ksys}
\teil Graph zeichnen: Zeichne den Graphen der Funktion f mit $f(x) = -4x + 3$ in das Koordinatensystem. \hfill (P10 2026 FOR)

\begin{ksys}[xmin=-4,xmax=4,ymin=-4,ymax=4] \end{ksys}
\end{teile}
\end{aufgabe}

%% TODO \verfahren{…}: Verfahrensüberschrift in Sachsprache fehlt in der Bank (2.3 b)
%% TODO Titel: Ich-kann-Satz fehlt in der Bank (Platzhalter: Kettenname) – Nr. 16
%% TODO Anweisung: ein Satz über der ersten Teilaufgabe fehlt in der Bank – Nr. 16
\begin{aufgabe}{Ablesen}
%% TODO \rechenplatz{n}: Schrittzahl des Grundfalls fehlt in der Bank (nur bei fester Schreibform, 2.3 b)
\begin{teile}
\teil n? Lies den y-Achsenabschnitt n der Geraden g ab.

\begin{ksys}[xmin=-5,xmax=5,ymin=-5,ymax=5,ablesen] \gerade{1}{3}{g} \end{ksys}
n = \leerfeld
\teil m? Lies die Steigung m der Geraden g mit einem Steigungsdreieck ab.

\begin{ksys}[xmin=-5,xmax=5,ymin=-5,ymax=5,ablesen] \gerade{2}{-3}{g} \end{ksys}
m = \leerfeld
\teil m? Lies die Steigung m der Geraden g mit einem Steigungsdreieck ab.

\begin{ksys}[xmin=-5,xmax=5,ymin=-5,ymax=5,ablesen] \gerade{-2}{1}{g} \end{ksys}
m = \leerfeld
\teil m? Lies die Steigung m der Geraden g ab. Gehe so weit nach rechts, bis du wieder auf einer Gitterlinie bist.

\begin{ksys}[xmin=-5,xmax=5,ymin=-5,ymax=5,ablesen] \gerade{0.5}{1}{g} \end{ksys}
m = \leerfeld
\teil Welche Gleichung gehört zur Geraden? Kreuze an. \\ \kreuz{$f(x) = 3x - 4$} \\ \kreuz{$f(x) = -4x + 3$} \\ \kreuz{$f(x) = 3x + 4$} \\ \kreuz{$f(x) = -3x - 4$}

\begin{ksys}[xmin=-5,xmax=5,ymin=-5,ymax=5,ablesen] \gerade{3}{-4}{g} \end{ksys}
\teil Welche Gleichung gehört zu welcher Geraden? Ordne den Geraden f, g und h je eine der Gleichungen zu: $y = -3x + 3$; $y = -\frac{1}{3}x + 3$; $y = -3$; $y = 2x + 3$; $y = -3x$; $y = 2x - 3$. \hfill (P10 2019 OS) \\

\begin{ksys}[xmin=-5,xmax=5,ymin=-5,ymax=5,ablesen] \gerade{-3}{3}{f} \gerade{2}{3}{g} \gerade{0}{-3}{h} \end{ksys}
f: \leerfeld \quad g: \leerfeld \quad h: \leerfeld
\end{teile}
\end{aufgabe}

%% TODO Titel: Ich-kann-Satz fehlt in der Bank (Platzhalter: Kettenname) – Nr. 17
%% TODO Anweisung: ein Satz über der ersten Teilaufgabe fehlt in der Bank – Nr. 17
\begin{aufgabe}{Parameter deuten (steigend, fallend, parallel, Sonderfall m = 0)}
\begin{teile}
\teil Welche zwei Geraden sind parallel? $f(x) = 4x - 1$, $g(x) = -4x + 5$, $h(x) = 4x + 6$
\end{teile}
\end{aufgabe}

%% TODO Titel: Ich-kann-Satz fehlt in der Bank (Platzhalter: Kettenname) – Nr. 18
%% TODO Anweisung: ein Satz über der ersten Teilaufgabe fehlt in der Bank – Nr. 18
\begin{aufgabe}{Graph zeichnen – Fehler finden, Begründen, Darstellungswechsel}
\begin{teile}
\teil Ich finde den Fehler: Paul liest am Steigungsdreieck ab: 1 nach rechts, 3 nach oben, also $m = \frac{1}{3}$. Benenne den Fehler und gib m richtig an.

\begin{ksys}[xmin=-4,xmax=4,ymin=-5,ymax=5,ablesen] \gerade{3}{-2}{g} \steigungsdreieck{0}{-2}{3} \end{ksys}
\teil Welche Gerade ist steiler: $f(x) = 3x - 7$ oder $g(x) = 5x + 4$? Begründe ohne Rechnung.
\teil Gleichung zur Geraden? Gib die Gleichung der Geraden g an.

\begin{ksys}[xmin=-5,xmax=5,ymin=-5,ymax=5,ablesen] \gerade{-1}{3}{g} \end{ksys}
f(x) = \leerfeld
\end{teile}
\end{aufgabe}
